\section{The Praxis of Trust}
\label{p13:sec:praxis}

\noindent
The paper arrives at its practical summit, and at the chapter toward which the whole has been moving. Everything established\,---\,the mechanism, the geometric reading, the criterion of justice, the structure of relational trust, the theory of translation, the optimisation of the juridical\,---\,has been knowledge \emph{for the sake of} practice, and the paper would betray its own deepest claim if it stopped at knowledge. For the series holds that to settle in the theoretical, to take the cognition of a thing as the completion of it, is itself an instance of the pathology the whole paper diagnoses\,---\,the failure of settlement, the symbol that declares a circuit closed which no one has walked. A paper on the generation of trust that ended in an account of the mechanism, without descending to the practice in which trust is actually generated and lived, would be such a closed and unwalked circuit. The praxis is therefore not an application appended to the theory; it is the theory's discharge of its own obligation not to settle. This chapter is long, and rises, at its centre, to the paper's metaphysical summit; the present part develops its ground, from the reason there must be a praxis at all to the honest reckoning with what no praxis can guarantee, and the part that follows it raises the summit and returns, at the last, to the family and the room.

\subsection{Why there must be a praxis}
\label{p13:sub:prax-why}

The reason the paper cannot end in theory is the paper's own. It has argued that the poetic symbol refuses settlement\,---\,that it presents generativity rather than a closed structure, and forces its meaning to be walked rather than cashed; and it has argued that the failure of settlement, the cashing of a sublimation no one has undergone, is the very form of the symbolic usurpation it diagnoses. To stop at the theoretical cognition of trust would be to commit this failure in the paper's own conduct: to cash, in the coin of understanding, a practice that has not been walked. The descent to praxis is thus demanded by the paper's content, and not merely by a wish to be useful. And the descent is also the place where the paper's argument and the author's life, which the Prelude held apart as the near and the far edge of one question, are to meet: the praxis of trust is what the author must do, before the family of the Prelude, and it is what the theory has been clearing the ground to make sayable. The chapter is the point at which the knowledge turns into a way of acting, and refuses, in turning, to settle.

\subsection{A polyphony of practices, and an ethical line across them}
\label{p13:sub:prax-polyphony}

There is no single practice of trust, and the series' polyphonic law holds in the practical register as it held in the theoretical: each of the frameworks the paper has drawn upon furnishes a practice of trust, and none governs the rest. The evidentialist practice builds the track record, accumulates the demonstrable, lets the facts mount until they speak. The deontological practice cultivates and shows a character, makes of the self a thing whose sincerity can be read. The relational practice tends the chains of \xref{p13:sec:relational}, invests in the couplings through which trust is transmitted. The Daoist practice, of which more below, declines strategy and lets the real show itself in its own time. The institutional practice works the endorsements and the registers. These are real and distinct practices, suited to different receivers and different moments, and the paper does not rank them; the practical wisdom of trust is in large part the wisdom of which practice a given situation calls for, and the polyphony is, here as throughout, not a defect to be resolved into a single method but the very shape of a competence adequate to an object no single framework possesses.

But across all these practices runs one ethical line, and on this the polyphony does not extend: descriptively plural, the practices are ethically subject, every one, to the single criterion of \xref{p13:sec:justice}. Whatever the practice\,---\,evidential, deontological, relational, Daoist, institutional\,---\,it is genuine where it produces an expectation by supplying a real ground that honours the other's free judgment and opens a relation that returns to him, and forged where it induces an expectation by manufacturing a signal that handles him as a mechanism and draws upon him as fuel. The plurality is a plurality of legitimate forms; the line between the genuine and the forged cuts across every one of them, and is the same line in each. The chapter is polyphonic in its description of the practices and monophonic in its ethic of them, and the next subsections turn first to a correction the polyphony itself requires, and then to the forged practice the ethic excludes.

\subsection{The justice of presentation}
\label{p13:sub:prax-presentation}

The polyphony of practices, as just set out, conceals a value-judgment that must be brought up and corrected, for it has quietly ranked the practices after all, and ranked them wrongly. There is a temptation, native to the Daoist and to a certain reading of the natural, to hold the practice of unstrategic self-showing\,---\,the letting of the real appear in its own time, without contrivance\,---\,as the highest and purest form of genuine trust-building, and to regard any \emph{active} presentation, any rendering of oneself in the form the receiver's cognition can best take, as already shading toward the manipulative. On this view, to present one's sincerity in the grammar of the evidentialist would be a falling-away from the purity of simply, sincerely being; and the active translation of \xref{p13:sec:translation} would carry a faint guilt, as a strategy where strategy ought not to be. This subsection reverses that judgment, and the reversal is the hinge of the whole praxis.

The reversal turns on a recognition the paper owes to its own theory of translation: that the active presentation of oneself in the grammar the other can read\,---\,the very thing the unstrategic ideal suspects\,---\,is not a fall from purity but, in the right conditions, the way of \emph{countering an epistemic injustice}. For epistemic injustice runs in two directions, and the unstrategic ideal sees only one of them. The first direction is the one everyone sees: the active evil of manipulation, the use of the other's grammar to manufacture a signal of something that is not there, which corrupts his epistemic autonomy (\xref{p13:sec:justice}). But there is a second direction the unstrategic ideal commits while believing itself innocent: the passive evil of \emph{refusing to translate}\,---\,of insisting on showing oneself only in one's own grammar and leaving the reception to fate. To say "I will simply appear, truly, in my own way, and the rest is not my concern" carries an unjust presupposition: that one's own grammar is the neutral and universal one, in which the real, shown plainly, will be legible to anyone of good faith. But there is no such neutral grammar (\xref{p13:sec:translation}); to show oneself only in one's own and require the other to read one there is to impose one's own grammar upon him as though it were the universal, and that imposition\,---\,however high-minded, however draped in the language of sincerity and refusal of contrivance\,---\,is itself an epistemic injustice, the more insidious for disguising itself as purity. The unstrategic ideal is not humble; it is, beneath the humility, the arrogance of treating one's own way of being legible as legibility itself.

The correction is therefore exact. Unstrategic self-showing is not the highest form of genuine trust; it is genuine only where the grammars happen to match, and where they do not\,---\,which is precisely the situation of the Prelude's suitor before the family\,---\,it is not purity but a concealed epistemic arrogance. And the active presentation of oneself in the other's grammar, which the unstrategic ideal suspected, is revealed as its opposite: not a strategy shading toward manipulation but an \emph{epistemic hospitality}, the actively going-toward the other's way of knowing, the taking-on of the labour of being legible to him in his terms rather than requiring him to learn mine. To present a real sincerity in the grammar the other can receive is to recognise his epistemology as legitimate and equal to one's own, and to move toward it; and this counters epistemic injustice rather than committing it, because it restores\,---\,rather than corrupts\,---\,the other's capacity to make, on his own epistemological terms, a true judgment about a real thing.

The line between the genuine and the forged is therefore not the line between the unstrategic and the active, between \emph{wu wei} and \emph{you wei}, between letting-be and doing. That was the error the unstrategic ideal made: it located the danger in activity as such. The line is the one it has always been: whether what is presented really exists. Genuine presentation may be unstrategic\,---\,the letting-appear, where the grammars match\,---\,or it may be active\,---\,the epistemic hospitality, the translation into the other's grammar, where they do not\,---\,and both are genuine, for both present, by whatever means, a thing that is real. Forged presentation is the active manufacture, in the other's grammar, of a signal for a thing that is not there. The difference is in the reality of what is presented, never in the activity or passivity of the presenting; and in the situation of epistemic mismatch\,---\,the situation that matters, the situation of the Prelude\,---\,the active presentation, the epistemic hospitality, is not the lesser and more suspect form but the one that more fully respects the other and more fully counters the injustice of requiring him to read one in a grammar not his own. To present one's real sincerity to the family in the form they can receive\,---\,the stable fact, the observable constancy, the evidentially legible\,---\,is not the abandonment of a deontological purity nor a step toward manipulation; it is respect for their epistemology, and the genuine active good, so long as what is presented is real.

\subsection{The forged practice excluded}
\label{p13:sub:prax-forged}

The ethic that runs across the polyphony excludes one practice absolutely, and it is worth naming the excluded practice precisely, for it is the practice that most resembles the genuine and is most apt to be mistaken for a legitimate technique. It is the practice of behavioural conditioning: the deliberate shaping of another's expectations by the engineering of reinforcements, the production of trust by the systematic management of what the other experiences as the consequence of trusting, so that the trust is trained into him as a conditioned response rather than grounded in any real self-commitment. It is the manipulation of \xref{p13:sec:justice} raised to a method, and its evil is threefold, the three aspects of the one wrong the justice criterion identified, here compounded.

It is, first, the treating of the other as a means and a mechanism\,---\,the Kantian wrong\,---\,for it works upon him as upon a system to be conditioned, bypassing the free judgment that genuine trust addresses. It is, second, the epistemic injustice\,---\,for it installs in him an expectation by training rather than by grounds, corrupting his capacity to form, on real grounds, a true judgment of his own. And it is, third and distinctively, a temporal forgery, a failure of settlement spread across time: it manufactures, by the management of reinforcements, the \emph{appearance of a constancy} that no real self-commitment underwrites, forging across time the very track record that the evidentialist reads as the shadow of a real sincerity\,---\,so that it is the precise isomorph, in the practice of trust, of the forged recommendation of \xref{p13:sec:relational}, which manufactures the symbolic form of a coupling that does not exist. The conditioned constancy is a decohered constancy: the observable shell of a track record, with no real coupling beneath it. The practice is excluded not because conditioning never produces reliable behaviour\,---\,it may\,---\,but because the reliability it produces is a forged precision, a trust trained into the other as into a mechanism, drawing upon him as fuel while corrupting his capacity to know, and forging, across the very dimension of time that should reveal the difference, the shadow that should have been the real's alone.

\subsection{Hume's gulf, and the honest reckoning with no guarantee}
\label{p13:sub:prax-hume}

Beneath the whole practice of trust lies a limit that no practice can cross, and the honesty of the praxis depends on reckoning with it rather than concealing it. It is Hume's gulf, met first in the diagnosis (\xref{p13:sub:diag-humekant}) and now to be faced in the practical register. No practice of trust, however genuine, can guarantee its own outcome. The one who builds the track record cannot guarantee he will keep faith tomorrow; the one who makes the costly self-commitment cannot guarantee the circumstance that will test it; the one who trusts cannot cross, by any quantity of evidence or any depth of self-binding, the gulf between what has been and what will be. The practitioner of trust is in exactly the position the prior paper of the series assigned to the practitioner of any good practice: he cannot guarantee that his practice is good, cannot occupy the position of an Other of infinite reason who could certify the absolute causal connection between his conduct and the good it aims at, cannot stand outside the gulf and warrant the crossing~\parencite{wanhong_grb,hume2000}. There is no view from which the keeping of the promise is assured in advance; there is only the promise, made and to be kept, across a gulf no guarantee spans.

This is not a counsel of despair but the precise location of the difference between the genuine and the forged at its deepest. For the difference between genuine and forged trust is not, at the last, a difference in the presence or absence of a guarantee\,---\,neither has a guarantee; the gulf is uncrossable for both. The difference is in the honesty with which the absence of a guarantee is borne. The forged practice pretends to a guarantee it does not have: it manufactures the appearance of an assured constancy, forging the certainty that the gulf forbids. The genuine practice makes its commitment in the full acknowledgement that nothing assures it\,---\,binds itself across the gulf precisely as one who knows the gulf is there and cannot be crossed by any warrant, and binds itself anyway. Genuineness is not the possession of a certainty the forged lacks; it is the honest bearing of an uncertainty the forged denies. The genuine promise is the one made in the knowledge that it cannot be guaranteed, and made, freely and without the false comfort of a guarantee, as a binding of oneself across a gulf one knows no evidence and no self-binding can close.

And here the praxis arrives at the threshold of its summit. For the question this raises is the one the rest of the chapter exists to answer: if no external guarantee can cross the gulf\,---\,if neither the sword, nor the evidence, nor the depth of the self-binding can assure that the promise will be kept\,---\,then what does the one who would genuinely bind himself \emph{do}, in the face of a gulf he cannot cross and a future he cannot warrant? The answer is not to find, at last, a guarantee that the chapter has overlooked; there is none. The answer is to make, across the gulf no guarantee spans, a free and self-grounding leap\,---\,to constitute, by an act of one's own freedom, the very constancy that no external warrant could supply. What that act is, how it is possible, and why it is the descent of something an older language would have called divine, is the matter of the summit to which the chapter now rises.

\subsection{Relational divinity: the descent of the Other into the self-legislating subject}
\label{p13:sub:prax-divinity}

\noindent
The leap the previous subsection named\,---\,the free constitution, across a gulf no warrant spans, of the very constancy no external guarantee could supply\,---\,is the paper's metaphysical summit, and it must be set out with the fullness its weight demands, for it is a claim the series has approached from many directions and never, until here, fully stated. It can be stated so. \emb{The self-legislation, self-adjudication, self-keeping, and self-removal of every recourse outside the law one has given oneself is the mechanism by which a relational divinity descends into the world.} This is the most original of the paper's claims, and it stands, the author believes, in no prior text of the series nor in the foundational paper; it is set out here in full and in emphasis as the thing toward which the whole has been moving.

Begin with the historical movement of the divine, which earlier work in the series established in its direction without specifying its present mechanism. There is a movement, across the history of how the divine has been located, in three stages. First, a \emph{transcendent} divinity: the god beyond the world, the divine power lodged in a heaven or a nature outside the human, before whom the oath is sworn and by whom its breach is avenged. Then, a \emph{symbolic} divinity: the divine descended into the Symbolic, into the great Other of the law, the institution, the symbolic authority, the impersonal reason before which the modern subject stands\,---\,the god replaced by the order of signs that judges and binds. And now, in the movement this paper completes, a \emph{relational} divinity: the divine descending out of the Symbolic and into the free practice of the relational subject, appearing no longer beyond the world nor in the order of signs above the subject, but in what the subject freely does, in relation, in the constituting of a bond.

The mechanism of this last descent\,---\,the present, concrete mechanism the earlier work marked the direction toward but did not supply\,---\,is the fourfold act. It is, first, \emph{self-legislation}, in the full Kantian sense of autonomy: the subject gives himself the law, lays it down as law for himself, so that the constancy he could not derive from the past (Hume's gulf forbids the derivation) is not a fact to be observed or predicted but a law to be \emph{enacted}, posited by the will, grounded in the \emph{ought} the subject lays upon himself\,---\,bypassing, rather than crossing, the gulf, by founding the constancy in an act of legislation rather than seeking it in an inference. It is, second, \emph{self-adjudication}, the act the author adds and on which the whole turns: the continuous holding of oneself to the law one has given, the standing as one's own judge over whether the law is kept, the self-arraignment and self-correction when one has strayed\,---\,which is the internalisation of Fuller's congruence (\xref{p13:sub:jur-fuller}), the bringing of one's conduct into agreement with one's announced rule, performed by the subject upon himself. A single free choice would be fragile, a resolution that the next moment might dissolve; but self-legislation joined to a continuous self-adjudication feeds each keeping of the law back into the loop that sustains it, and the constancy is thereby \emph{generated}\,---\,turned out, moment by moment, by the self-legislating and self-adjudicating loop\,---\,rather than guaranteed in advance. It is, third, \emph{self-keeping}, the actual holding to the law across time. And it is, fourth, the \emph{self-removal of every recourse outside the law one has given}\,---\,which is the deepest sense of the burning of the bridge, met first in \xref{p13:sub:mech-likelihood} and now to be completed.

For here the forward reference of the mechanism is discharged. The burning of the bridge was described, at \xref{p13:sub:mech-likelihood}, in its aspect as a high-fidelity signal\,---\,something done to be read by another, furnishing his model with hard-to-forge evidence. That was, it was said there, only one of its two faces and not the deeper. The deeper face is this: that one forecloses one's own recourse not in order to send a message but in order to constitute, for oneself, a law\,---\,\zh{我立此为法}, I have laid this down as law, and so that road no longer exists for me. The external law closes the road by force, mechanically, by the threat at the road's end; self-legislation closes the road by freedom, by the laying-down that makes the road, for this subject, not a road at all. The reader who reaches this point will recognise that the trunk was pointing here all along: that beneath the burning of the bridge as a signal lay the burning of the bridge as self-legislation, and that the costly act read by another was, in its depth, a subject's free and irrevocable binding of himself, in the manner in which a god is said to bind himself in a covenant\,---\,leaving no recourse, because the divine does not keep a recourse. The self-removal of recourse is the point at which the human exceeds the ordinary human, in which the part of the human capable of self-transcendence, capable of binding itself irrevocably as a god binds himself in a covenant, is realised; and it is the divine core of the vow.

\subsubsection{Why the divinity is relational and not a solitary self-deification}
\label{p13:sub:prax-relational-not-solitary}

The gravest danger of this claim must be met head-on, for a self-legislation that made itself its own law and its own judge would seem the very figure of hubris\,---\,the subject deifying himself, circling his own image, the solitary self-grounding that is the oldest form of the imaginary's idolatry. The answer to the danger is the single consideration on which the whole claim's safety rests, and it is ontological. \emb{The subject is not first a self and then a party to relations; the subject is constituted by relation, is relational through and through, before and beneath any self it could deify.} On the relational ontology the series has held, there is no prior, separate self that could turn upon itself in a closed circle of self-worship; the self that legislates is already, in its very constitution, relational\,---\,made of the relations it stands in. And therefore the law it gives itself, and the divinity that descends in the giving, are relational from the first. The self-legislation is never "I, to myself" in a closed loop; it is "I, as a relational being, within the relations that constitute me, for those relations, toward those relations"\,---\,the laying-down, by a self that is made of its bond to the other, of a law for the keeping of that bond. The betrothal vow is exactly this: a self-legislation performed by a self constituted in the relation to the one it vows to, laying down, for that relation and within it, the law of how it will be kept.

This relational constitution is, further, what dissolves the tension that would otherwise destroy self-adjudication. A solitary self-judge is the surest route to self-deception: nothing is easier than to acquit oneself, to flatter the law into agreement with the conduct, when one is one's own only judge. But the relational self-adjudication is not solitary\,---\,it is conducted, of necessity, within the regard of the relation. The other's witness, the gaze that the empty place of the \emph{das Ding} holds open, is not an external check added to the self-adjudication from outside; it is internal to the very structure of being a relational subject, which is constituted in being-seen as much as in seeing. The honesty of the self-adjudication is therefore not secured by an external sword\,---\,there is none, and the gulf forbids one\,---\,but by the ontological fact that the self which adjudicates is a self already constituted in relation, already under a regard it did not add and cannot remove. Relational divinity carries its own guard against self-deception, and the guard is not a borrowed sword but the subject's own relational being.

\subsubsection{The personal-divine form of positive holonomy}
\label{p13:sub:prax-holonomy-divine}

This fourfold act is the personal and divine form of the positive, adiabatic holonomy of \xref{p13:sec:holonomy}. The loop closes\,---\,the constancy is sustained, turn upon turn\,---\,not because an external pump drives it (no transcendent god, no symbolic Other, no sword: these would be the non-adiabatic driving, the borrowed phase), but because the relational subject's own self-legislating and self-adjudicating loop closes it from within, endogenously, accumulating its phase as its own. It is the discharge, at the level of the person, of the claim of \xref{p13:sec:freedom}, that autonomy is the highest form of freedom\,---\,not the merely negative freedom of being undetermined, but the positive freedom of giving oneself the law. And it gathers to itself the deepest voices of the series and of the traditions it has drawn upon: the foundational paper's intimation that the tendency of the Other toward its final dissolution opens the way to a renewed approach to a pure relational being~\parencite{wanhong_grb}; the Buddhist teaching that mind is itself the Buddha, that all dharmas are made by mind\,---\,\zh{是心是佛，一切法由心造}; the Daoist refusal of the external warrant; and the betrothal document's own closing line, \zh{本乎此心，自成永恒}\,---\,that the eternal is founded not upon any external guarantee but here, in this heart, which is itself the ground.

\subsubsection{The inversion of value: the non-binding vow as the more divine}
\label{p13:sub:prax-inversion}

A consequence follows that inverts the ordinary ranking of instruments, and it is the point at which the whole paper's central paradox is resolved into a claim. The ordinary ranking holds the marriage certificate above the betrothal vow, the contract above the memorandum of understanding, because the former in each pair carries the backing of the state, the force of law, the reliability of the sword. From the standpoint of relational divinity, the ranking is exactly inverted. \emb{The memorandum of understanding is more divine than the contract, the betrothal vow more divine than the marriage certificate}\,---\,precisely because they carry no force, have no external guarantee to borrow, and so can be sustained only by the subject's own self-legislation, self-adjudication, and self-keeping in the real.

The logic is this. The contract and the marriage certificate \emph{outsource} the threefold function\,---\,legislation, adjudication, keeping\,---\,to the symbolic Other: the law is the state's, the judging is the court's, the keeping is compelled by the sword. The parties borrow an external, symbolic divinity to underwrite their bond, and need not legislate, adjudicate, or foreclose their own recourse, for the functions are outsourced. The memorandum and the vow, having no force, force these functions back upon the subject: with no external sword, the legislation can only be self-legislation; with no external court, the adjudication can only be self-adjudication; with no external compulsion, the keeping can only be self-keeping and the self-removal of recourse. The parties have nowhere to borrow an external divinity, and must therefore \emph{become} the legislator, the judge, and the keeper themselves\,---\,must let the relational divinity descend into themselves, or the bond does not hold. The "defect" of having no force is thus the very condition of the divinity: precisely because no external sword keeps the vow, the keeping must issue from one's own free, irrevocable self-binding, which is the realisation, and not the absence, of the divine. This completes the geometric reading of \xref{p13:sub:holo-convergence}: the contract and the certificate are non-adiabatic, their loop held closed by the external pump of the state's force, their phase borrowed; the memorandum and the vow are adiabatic, their loop closed by nothing but the subject's endogenous self-legislation, their phase their own\,---\,the purer specimen of genuine holonomy, into which the certificate admits an admixture of the external pump. And it raises the counter-intuitive claim of the diagnosis to its summit: the betrothal vow is not merely, as was first said, possibly more productive of real trust than the complete contract; it is, in the dimension of the divine, the purer thing, because it forces the divine back into the free practice of the subject rather than borrowing it from an Other.

The honest boundary must be drawn, lest the inversion be pressed past what it claims. It is not said that the memorandum and the vow are superior to the contract and the certificate in every dimension. In the real functions of protecting the vulnerable, deterring betrayal, and providing a remedy, external force remains necessary and is not to be romantically discarded\,---\,exactly as \xref{p13:sec:judicial} insisted in relocating rather than abolishing the sword. The claim is narrow and precise: that in the dimension of the divine, of the genuine, the non-coercive instrument is the purer, because it forces the divinity back into the subject's own free practice. The weak still need the sword's protection at the scale of institutions; but the divinity of a genuine bond was never lodged in the sword, and is found, at the scale of relation, only where the sword is not. It is not, therefore, that one should forgo the certificate for the vow; it is that, beyond and above and before the certificate, what makes the bond a bond at all is the self-legislating relational divinity, and that this is most purely realised in the vow that has no force.

\subsection{The concrete operation}
\label{p13:sub:prax-operation}

From the summit the chapter descends to the concrete, for the praxis must issue in something one can actually do, and the abstractions of divinity must be cashed\,---\,here legitimately, for here the path is walked and not merely declared\,---\,into an operation. The operation has three movements, and a condition of time.

The first movement is to \emph{identify the other's grammar of trust}, and its configuration (\xref{p13:sub:trans-configuration}): to discern, of the one whose trust is to be won, which of the grammars of \xref{p13:sub:trans-grammars} dominates his reading, in what weighting, so that one knows the language into which one's sincerity must be translated. One does not present in the abstract; one presents to this configuration, and the discernment of it is the first and indispensable act.

The second movement is to \emph{produce, in the other's grammar, a real possibility of expectation}\,---\,to translate, by the paths of \xref{p13:sec:translation}, and above all by the epistemic hospitality of \xref{p13:sub:prax-presentation}, one's real sincerity into the form the other's configuration can receive. For the evidentialist family, this is the projection of the real into the register of the observable and the accomplished: not the manufacture of a track record one does not have, which would be forgery, but the patient letting of one's real self-legislated constancy cast, in the evidential register, the shadow by which it can be recognised. The thing presented must be real\,---\,must be a sincerity that exists and a constancy that is genuinely self-legislated\,---\,and it must be presented in the grammar the other reads. This is the whole of the second movement: a real thing, hospitably translated.

The third movement is the \emph{right use of relational trust} (\xref{p13:sec:relational}), and it carries the chapter's most delicate ethical discrimination. The trust of the family runs, very largely, along the chain that connects the suitor to them through the one they already love\,---\,she is, in the structure of the network, his recommender, and the trust may reach them along that link as it reaches along no other. The right use of this is to let the trust transmit because the coupling it transmits is real\,---\,because she is, in the genuine sense, the co-author of the bond, and her steady bearing toward it is the unfakeable real component (\xref{p13:sub:rel-threeregisters}) that the family, even reading evidentially, can read. The wrong use\,---\,the forged use, excluded by the justice criterion\,---\,is to instrumentalise her as a mere channel, to deploy her to "work on" the family, to siphon her relational capital toward a coupling between the suitor and the family that has not in fact occurred. The line is exactly the line of \xref{p13:sec:justice}: between a real coupling that transmits because it is real, and a relational capital siphoned to transmit a coupling that is not there. The one honours her as co-author; the other renders her fuel.

And the condition of time. None of this can be hurried, and the attempt to hurry it is itself a tell of the forged. The self-legislated constancy casts its evidential shadow only as it is actually lived across time; the real coupling transmits only as it is really borne; the shared grammar of mutual translation grows only in the continued interaction. Time cannot manufacture the coupling-trust, which is real from the first or not at all; but time is what lets the real cast, in the register the evidentialist reads, the shadow sufficient to be recognised. The praxis is therefore not a technique that produces trust on demand but a way of living that lets a real thing become, across time, recognisable to one whose grammar is not one's own. It cannot be accelerated; it can only be deepened.

\subsection{The return to the room}
\label{p13:sub:prax-room}

The chapter, and with it the body of the paper, returns at the last to the scene from which the Prelude set out: to the one who must sit across from a family not yet his own, and be trusted. Everything the paper has built converges on that room, and it converges not on a method for controlling what happens there but on a way of being present in it.

For the end of the praxis is not the mastery of the other's trust\,---\,trust cannot be mastered, only generated, and generated across a gulf no one can guarantee. The end is to be ready to be truly present: to have become one who has a real sincerity to present, and to present it in the grammar the family can receive, with the epistemic hospitality that honours their way of knowing rather than demanding they learn his. What he can give them is not an external guarantee\,---\,he has none to give, and to manufacture the appearance of one would be the forgery the whole paper has worked to refuse. What he can give them is the sight of a person who legislates his own law and holds himself to it\,---\,who has laid down, for the relation that constitutes him, the law of how he will keep faith, and who adjudicates himself against it continually, in the open, where it can be seen. This constancy is real, for it is the product of the self-legislating loop; and being real, it casts, in time and in the evidential register the family reads, a true shadow\,---\,a track of consistent, self-correcting keeping that the evidentialist can read as what it is, the unforgeable sign not of a compelled obedience nor of a performed display but of a free subject keeping the law he gave himself. That is what there is to give: not a guarantee, which cannot be given, but the presence of a self-legislating relational subject, whose constancy is its own and whose keeping of faith is the descent, in him, of a divinity that needs no sword\,---\,offered, across the gulf, to be trusted or not, freely, by those he cannot compel and would not if he could.
